% Options for packages loaded elsewhere
\PassOptionsToPackage{unicode}{hyperref}
\PassOptionsToPackage{hyphens}{url}
\PassOptionsToPackage{dvipsnames,svgnames,x11names}{xcolor}
%
\documentclass[
  10pt,
  a4paper,
]{article}
\usepackage{amsmath,amssymb}
\usepackage{setspace}
\usepackage{iftex}
\ifPDFTeX
  \usepackage[T1]{fontenc}
  \usepackage[utf8]{inputenc}
  \usepackage{textcomp} % provide euro and other symbols
\else % if luatex or xetex
  \usepackage{unicode-math} % this also loads fontspec
  \defaultfontfeatures{Scale=MatchLowercase}
  \defaultfontfeatures[\rmfamily]{Ligatures=TeX,Scale=1}
\fi
\usepackage{lmodern}
\ifPDFTeX\else
  % xetex/luatex font selection
  \setmainfont[]{Liberation Serif}
  \setsansfont[]{Liberation Sans}
  \setmonofont[]{DejaVu Sans Mono}
\fi
% Use upquote if available, for straight quotes in verbatim environments
\IfFileExists{upquote.sty}{\usepackage{upquote}}{}
\IfFileExists{microtype.sty}{% use microtype if available
  \usepackage[]{microtype}
  \UseMicrotypeSet[protrusion]{basicmath} % disable protrusion for tt fonts
}{}
\makeatletter
\@ifundefined{KOMAClassName}{% if non-KOMA class
  \IfFileExists{parskip.sty}{%
    \usepackage{parskip}
  }{% else
    \setlength{\parindent}{0pt}
    \setlength{\parskip}{6pt plus 2pt minus 1pt}}
}{% if KOMA class
  \KOMAoptions{parskip=half}}
\makeatother
\usepackage{xcolor}
\usepackage[margin=22mm]{geometry}
\setlength{\emergencystretch}{3em} % prevent overfull lines
\providecommand{\tightlist}{%
  \setlength{\itemsep}{0pt}\setlength{\parskip}{0pt}}
\setcounter{secnumdepth}{-\maxdimen} % remove section numbering
\ifLuaTeX
\usepackage[bidi=basic]{babel}
\else
\usepackage[bidi=default]{babel}
\fi
\babelprovide[main,import]{korean}
\ifPDFTeX
\else
\babelfont{rm}[]{Liberation Serif}
\fi
% get rid of language-specific shorthands (see #6817):
\let\LanguageShortHands\languageshorthands
\def\languageshorthands#1{}
\usepackage{xeCJK}
\setCJKmainfont{Noto Serif CJK KR}
\setCJKsansfont{Noto Sans CJK KR}
\setCJKmonofont{Noto Sans Mono CJK KR}
\xeCJKsetup{CJKspace=true}
\usepackage{xurl}
\usepackage{fvextra}
\usepackage{fancyhdr}
\AtBeginDocument{\DefineVerbatimEnvironment{verbatim}{Verbatim}{breaklines=true,breakanywhere=true,fontsize=\small}}
\setlength{\emergencystretch}{4em}
\setlength{\headheight}{15pt}
\pagestyle{fancy}\fancyhf{}
\fancyhead[L]{\small 중성자별 구조와 비방사 진동}
\fancyfoot[C]{\thepage}
\renewcommand{\headrulewidth}{0.3pt}
\setcounter{tocdepth}{2}
\ifLuaTeX
  \usepackage{selnolig}  % disable illegal ligatures
\fi
\usepackage{bookmark}
\IfFileExists{xurl.sty}{\usepackage{xurl}}{} % add URL line breaks if available
\urlstyle{same}
\hypersetup{
  pdftitle={A01 · 단위·밀도·주파수의 경계},
  pdflang={ko-KR},
  colorlinks=true,
  linkcolor={Maroon},
  filecolor={Maroon},
  citecolor={Blue},
  urlcolor={Blue},
  pdfcreator={LaTeX via pandoc}}

\title{A01 · 단위·밀도·주파수의 경계}
\usepackage{etoolbox}
\makeatletter
\providecommand{\subtitle}[1]{% add subtitle to \maketitle
  \apptocmd{\@title}{\par {\large #1 \par}}{}{}
}
\makeatother
\subtitle{과제 큐시트}
\author{}
\date{}

\begin{document}
\maketitle

\setstretch{1.12}
\textbf{선행:} A00　 \textbf{독립 작업:} 1--2시간 (계획치)

\subsection{왕복 변환이 맞으면 단위도 맞은
것일까?}\label{uxc655uxbcf5-uxbcc0uxd658uxc774-uxb9deuxc73cuxba74-uxb2e8uxc704uxb3c4-uxb9deuxc740-uxac83uxc77cuxae4c}

순변환의 계수를 잘못 정하고 역변환에 그 역수를 썼다고 하자. 왕복 시험과
독립 기준값 시험 중 어느 쪽이 이 오류를 찾을지 먼저 선택하자.

먼저 한두 문장으로 답한 뒤 아래 설명을 읽자. 첫 예측이 틀렸거나 아직
모르겠다고 적는 것은 감점하지 않는다.

\subsection{필요한 관계와 읽을
위치}\label{uxd544uxc694uxd55c-uxad00uxacc4uxc640-uxc77duxc744-uxc704uxce58}

\(1\,\mathrm{MeV\,fm^{-3}}=1.602176634\times10^{32}\,\mathrm{J\,m^{-3}}\),
\(\rho_0=m_b n\), \(\rho_E=\epsilon/c^2\), \(f=\omega/(2\pi)\).

\begin{itemize}
\tightlist
\item
  \path|notes/chapters/ch02_units.md|
\item
  \path|notes/equations/conventions.md|
\item
  \path|notes/chapters/ch02_units.md|
\end{itemize}

\subsection{계산으로
옮기기}\label{uxacc4uxc0b0uxc73cuxb85c-uxc62euxae30uxae30}

상수와 단위 모듈은 제공한다. 스스로 남길 것은 정의에서 얻은 기준값과 그
계산의 이유다.

\begin{enumerate}
\def\labelenumi{\arabic{enumi}.}
\tightlist
\item
  환산의 분모 fm 세 개를 손으로 전개한다.
\item
  제공 단위 모듈을 사용하여 기준값과 왕복을 각각 검사한다.
\item
  같은 잘못된 계수를 쓴 순·역변환이 왕복을 통과하는 반례를 만든다.
\item
  내부에너지가 있는 상태에서 두 질량밀도의 차이를 해석한다.
\end{enumerate}

\newpage

\subsection{실행해서 무엇을 볼
것인가}\label{uxc2e4uxd589uxd574uxc11c-uxbb34uxc5c7uxc744-uxbcfc-uxac83uxc778uxac00}

패키지 루트에서 실행한다. \texttt{-\/-plan}은 경로만 확인하며 계산
성공을 뜻하지 않는다.

\begin{verbatim}
python tools/run_assignment.py A01 --plan
python tools/run_assignment.py A01 --run \
  --workspace learner --output ../A01_result.json
\end{verbatim}

\subsection{결과를
설명하기}\label{uxacb0uxacfcuxb97c-uxc124uxba85uxd558uxae30}

왕복 성공은 두 연산의 일관성이다. 올바른 물리 단위인지는 정의에서 구한
한 값으로 따로 확인해야 한다. 두 밀도와 두 주파수의 이름도 이때
구별한다.

\textbf{남길 증거.} 환산 유도와 차원표; 독립 기준값과 왕복 결과; 2π
누락을 차원만으로 찾을 수 없는 이유

\textbf{판단 기준.} 2점 예: 왕복만 맞는 틀린 환산을 제시하고 독립
anchor가 이를 배제함을 정량적으로 보여준다. 최초 예측의 정오보다, 결과를
본 뒤 이유를 고치고 다른 입력에도 적용하는지를 본다.

\subsection{막혔을 때}\label{uxb9c9uxd614uxc744-uxb54c}

\textbf{힌트 1.} 변환 전 이름과 단위를 열 제목에 함께 적는다.

\textbf{힌트 2.} 왕복과 독립 anchor를 다른 주장으로 기록한다.

\subsection{다음 만남에 다시
묻기}\label{uxb2e4uxc74c-uxb9ccuxb0a8uxc5d0-uxb2e4uxc2dc-uxbb3buxae30}

다음 회차에는 환산표를 덮고 분모의 fm 세 개만 다시 전개한다. 그 뒤 표를
열어 지수와 단위를 고친다. 먼저 짧게 답한 뒤 본문과 자신의 기록을 열어
확인한다.

\textbf{범위.} 단위시험은 EOS의 물리성이나 별의 질량을 검증하지 않는다.

\end{document}
